\section{Methodology}\label{sec:method}

\subsection{System overview}
Agentic Tree-RAG separates source organization, query planning, execution, and answer evidence. At ingestion time, Data Agents profile heterogeneous files and create a hierarchical Description Tree. At query time, a planner creates an executable Query Tree over that source index. NodeAgents retrieve candidate source descriptions, the MainAgent invokes registered tools, and JEV makes evidence-conditioned tree-operation decisions. QueryTreeManager is the sole component that applies those decisions and updates execution state. A distinct Evidence Tree records retrieved source pointers, tool outputs, and decision traces for final synthesis.

\subsection{Joint source descriptions and geographic ontology}
For each ingested source, the Data Agent extracts an observable profile, such as file and field metadata, value samples, geometry types, and coordinate information when available. The LLM jointly generates (i) a retrieval-oriented description and (ii) a proposed, source-grounded domain ontology. The ontology contains candidate classes, subclass relations, properties, field links, and bindings to the source description. Structural validation checks proposed terms against the supplied profile; the ontology is retained with source provenance and is not treated as verified instance data. Description nodes store semantic text, lexical representations, and pointers to the original records. This keeps the ontology aligned with the same source and ingestion context as its description tree.

\subsection{Hybrid source retrieval}
For each retrieval task, the system searches the applicable Description Tree scope using BM25 over lexical representations and modality-specific HNSW over embeddings. Reciprocal-rank fusion combines the candidate lists, and a cross-encoder reranks the fused candidates. Retrieval identifies relevant sources and record pointers; it does not by itself establish a geographic fact or produce the final answer. The MainAgent uses the retrieved pointers and task specification to select an available table, geometry, document, vision, or GeoSPARQL tool.

\subsection{Query Tree representation and planning}
The planner represents an investigation as a rooted tree of retrieval, computation, and synthesis tasks. Each node specifies its query, target, modality, expected tool, evidence goal, input producers, execution policy, and declared outcome routes. Parent-child edges express control flow; input references express data dependencies, so a task can consume evidence from a different branch without conflating dependency with branch activation. An \texttt{always} child runs for every parent outcome when its inputs are ready. A \texttt{conditional} child runs only when its activation condition and declared parent route match. Success, unresolved, and other supported outcomes remain distinct; missing or incomplete evidence is not silently interpreted as a negative result. The planner validates route references and the combined control/dataflow graph before execution, and constructs multi-stage paths only where the requested operation requires them.

\subsection{Asynchronous execution and adaptive decisions}
The Orchestrator dispatches independent ready branches through a bounded worker pool. A NodeAgent retrieves candidates for one task, while the MainAgent is the sole caller of registered tools. JEV receives the task, result status, and available evidence-grounded follow-up candidates, then chooses a permitted tree operation such as completion, pruning of an irrelevant conditional branch, or adding a supplied follow-up. It does not rewrite source evidence or invent a new candidate query. QueryTreeManager validates each operation, enforces dependencies and execution policy, and alone mutates the Query Tree. Results from completed, unresolved, and pruned paths remain traceable. This design permits concurrent independent work while preserving ordered execution where a task depends on upstream evidence.

\subsection{Spatial tools and GeoSPARQL}
Spatial processing uses the tool appropriate to the data and declared task. Tabular and geometry tools operate on retrieved source records and their observed fields and coordinates. When ingestion has established a source binding to an ontology-backed graph with usable geometry, the planner can route a spatial relation or measurement to GeoSPARQL. The query uses the mapped geographic vocabulary and the graph's verified geometry binding; its returned bindings are recorded as tool evidence. GeoSPARQL supplements source retrieval and ordinary data tools: it does not replace the Description Tree, infer unsupported category membership, or turn an unverified ontology proposal into a fact. If no executable binding exists, the plan must use another supported evidence path or preserve the limitation as unresolved.

\subsection{Evidence-grounded synthesis}
Each task contributes source identifiers, retrieved records, tool results, statuses, and execution decisions to the Evidence Tree. The final synthesis consumes the outputs declared by its inputs, distinguishes measured results from semantic/category evidence, and retains source attribution and unresolved alternatives. Consequently, each answer component can be traced through the execution plan to the records or tool outputs that support it.
